\documentclass[a4paper,12pt]{article}
\usepackage[utf8]{inputenc}
\usepackage[T1]{fontenc}
\usepackage{lmodern}
\usepackage{geometry}
\usepackage{setspace}
\usepackage{fancyhdr}
\usepackage{hyperref}
\usepackage{amsmath}
\usepackage{amsfonts}
\usepackage{amssymb}
\usepackage{mathtools}
\usepackage{float}
\usepackage{tabularx}
\usepackage{enumitem}
\usepackage{graphicx}
\usepackage{ragged2e}
\usepackage[spanish]{babel}
\usepackage{booktabs}
\usepackage{longtable}
\usepackage{multirow}
\usepackage{xcolor}
\usepackage{microtype}
\usepackage{caption}
\usepackage{textcomp}
\newcommand{\euro}{\texteuro}
\newcommand{\E}{\mathbb{E}}
\newcommand{\Var}{\mathrm{Var}}
\newcommand{\Cov}{\mathrm{Cov}}
\newcommand{\VaR}{\mathrm{VaR}}
\newcommand{\TVaR}{\mathrm{TVaR}}
\addto\captionsspanish{\renewcommand{\tablename}{Tabla}}
\addto\captionsspanish{\renewcommand{\figurename}{Figura}}
\captionsetup{font=small,labelfont=bf}

% Configuración de márgenes
\geometry{left=2.5cm,right=2.5cm,top=3cm,bottom=3cm}

% Configuración de encabezados y pies de página
\pagestyle{fancy}
\fancyhf{}
\renewcommand{\headrulewidth}{0pt}
\fancyfoot[C]{\thepage}

% Configuración de interlineado
\onehalfspacing
\justifying

% Configuración de hipervínculos
\hypersetup{
    hidelinks,
    colorlinks=false
}

\begin{document}

% Primera página (Título, Resumen y Abstract)
\begin{titlepage}
    \centering
    \vspace*{0.1cm}
    \huge\textbf{Tema 3: Preguntas de repaso}\\
    \vspace{1cm}
    \small{Luiz Alberto Luzuriaga Izquierdo}\\
    \vspace{0.1cm}
    Facultad de ciencias económicas, empresariales y turismo. \\
    \vspace{0.1cm}
    Universidad de Alcalá de Henares. \\
    \vspace{0.1cm}
    Curso 2026-2027\\
    \vspace{0.5cm}

{\bfseries Resumen}

\vspace{0.1cm}
\noindent
\begin{justify}
Respuestas razonadas a las tres preguntas de repaso del tema. Tema 3 (Modelo colectivo de riesgo): Frecuencia y severidad, recursión de Panjer y FFT, sobredispersión y colas pesadas.\\\\
\vspace{0.1cm}
\noindent
\textbf{Palabras clave}: modelo colectivo, frecuencia, severidad, Panjer, preguntas de repaso.
\end{justify}
\begin{center}
    {\bfseries Abstract}
\end{center}
\vspace{0.1cm}
\noindent
\begin{justify}
Class work for Topic 3 of Actuarial Mathematics of Stability and Solvency of Insurance Activity: review questions with actuarial reasoning and numerical evidence.\\\\
\vspace{0.1cm}
\noindent
\textbf{Keywords}: collective risk model, frequency, severity, Panjer, review questions.
\end{justify}

\end{titlepage}

% Índice
\newpage\pagenumbering{gobble}
\renewcommand{\contentsname}{Índice}
\newgeometry{left=4cm, right=4cm, top=3cm, bottom=2.5cm}
\tableofcontents

\restoregeometry
\newpage
\numberwithin{equation}{section}
\pagenumbering{arabic}
\setcounter{page}{1}

\section{Tema 3: Modelo colectivo de riesgo}
\label{sec:tema3repaso}
\noindent
Respuestas razonadas a las tres preguntas de repaso del tema. Material de referencia: diapositivas del Tema 3 (Prof. Hugo González Riera).

\subsection{Pregunta 1 -- Por qu\'e $\Var(S)$ tiene dos t\'erminos}

\noindent\textbf{Enunciado.} Explica con tus propias palabras por qu\'e $\Var(S)$ tiene dos t\'erminos en el modelo colectivo ($\E[N]\Var(X)$ y $\Var(N)\E[X]^2$) y qu\'e mide cada uno.

\subsubsection*{Derivaci\'on}
$S=\sum_{i=1}^NX_i$ con $N$ aleatorio, $X_i$ i.i.d. e independientes de $N$. La esperanza condicionada es $\E[S\mid N]=N\,\E[X]$ y la varianza condicionada $\Var(S\mid N)=N\,\Var(X)$ (suma de $N$ variables independientes). La f\'ormula de la varianza total da
\[
\Var(S)=\underbrace{\E\big[\Var(S\mid N)\big]}_{\E[N]\,\Var(X)}+\underbrace{\Var\big(\E[S\mid N]\big)}_{\Var(N)\,\E[X]^2},\qquad \E[S]=\E[N]\,\E[X].
\]
(F\'ormula de Wald; ambas pueden deducirse tambi\'en derivando $M_S(t)=M_N(\ln M_X(t))$.)

\subsubsection*{Qu\'e mide cada t\'ermino}
\begin{itemize}[leftmargin=*]
\item $\E[N]\Var(X)$ -- \textbf{riesgo de severidad (o de proceso)}: aun sabiendo exactamente cu\'antos siniestros habr\'a, el coste total fluct\'ua porque cada uno cuesta una cantidad distinta. Es la media de la varianza \emph{dentro} de cada n\'umero de siniestros.
\item $\Var(N)\E[X]^2$ -- \textbf{riesgo de frecuencia}: aun si todos los siniestros costaran exactamente $\E[X]$, el total fluct\'ua porque cambia cu\'antos hay. Es la varianza de la media condicionada: se transmite al total escalada por el cuadrado del coste medio.
\end{itemize}

\subsubsection*{Con n\'umeros}
\begin{center}
\begin{tabular}{lrrrr}
\toprule
Cartera & $\E[N]\Var(X)$ & $\Var(N)\E[X]^2$ & $\Var(S)$ & peso de la frecuencia\\
\midrule
Auto (Poisson 1.200; Gamma(2; 450)) & $4{,}86\cdot10^{8}$ & $9{,}72\cdot10^{8}$ & $1{,}458\cdot10^{9}$ & $66{,}7\%$\\
Hogar (BN(50; 3); Pareto(4; 3.000)) & $3{,}00\cdot10^{8}$ & $6{,}00\cdot10^{8}$ & $9{,}00\cdot10^{8}$ & $66{,}7\%$\\
\bottomrule
\end{tabular}
\end{center}
Aunque en la intuici\'on se tienda a pensar que el riesgo est\'a en la severidad, en ambas carteras dos tercios de la varianza vienen de la frecuencia. Con Poisson, $\Var(N)=\E[N]$ y $\Var(S)=\E[N]\,\E[X^2]$, de modo que el peso de la frecuencia es $\E[X]^2/\E[X^2]=1/(1+CV_X^2)$: con $CV_X^2=1/\alpha=0{,}5$ resulta $2/3$.

\subsubsection*{Lo que ense\~na la descomposici\'on}
La variaci\'on relativa es
\[
CV_S^2=\frac{\Var(S)}{\E[S]^2}=\frac{\Var(N)}{\E[N]^2}+\frac{CV_X^2}{\E[N]} .
\]
El segundo sumando se diversifica al crecer la cartera ($\propto1/\E[N]$). El primero no, si $\Var(N)/\E[N]^2$ no tiende a cero: con Poisson vale $1/\lambda$ y se anula, pero con Binomial Negativa vale $1/r$ y queda como un \emph{suelo}. En Hogar, $1/r=0{,}02$ implica $CV_S\ge\sqrt{0{,}02}=14{,}1\%$ por grande que sea la cartera (el valor total es $CV_S=20\%$). En Auto con Poisson, $CV_S=3{,}5\%$ y baja como $1/\sqrt\lambda$.

\paragraph{Conclusi\'on.} Los dos t\'erminos son los dos or\'igenes del riesgo (cu\'anto cuesta cada siniestro y cu\'antos hay). Distinguirlos importa porque el de severidad se diversifica con el volumen y el de frecuencia, si hay sobredispersi\'on, no del todo.

\bigskip\hrule\bigskip

\subsection{Pregunta 2 -- La Binomial Negativa como mezcla Poisson--Gamma}

\noindent\textbf{Enunciado.} ?`Por qu\'e la Binomial Negativa se interpreta como una mezcla Poisson--Gamma y qu\'e papel juega en la sobredispersi\'on observada en carteras reales?

\subsubsection*{Demostraci\'on}
Sea $\Lambda\sim\text{Gamma}(r,\text{escala }\beta)$, con densidad $\lambda^{r-1}e^{-\lambda/\beta}/(\Gamma(r)\beta^r)$, y $N\mid\Lambda=\lambda\sim\text{Poisson}(\lambda)$. Entonces
\[
P(N=n)=\int_0^\infty\frac{e^{-\lambda}\lambda^n}{n!}\,\frac{\lambda^{r-1}e^{-\lambda/\beta}}{\Gamma(r)\beta^r}\,d\lambda=\frac{\Gamma(n+r)}{n!\,\Gamma(r)}\Big(\frac1{1+\beta}\Big)^{r}\Big(\frac{\beta}{1+\beta}\Big)^{n},
\]
que es exactamente la Binomial Negativa de las diapositivas. De forma equivalente, la funci\'on generatriz de probabilidades es $P_N(z)=\E\big[e^{\Lambda(z-1)}\big]=\big(1-\beta(z-1)\big)^{-r}$.

\subsubsection*{Momentos y sobredispersi\'on}
Por la ley de la varianza total,
\[
\E[N]=\E[\Lambda]=r\beta,\qquad \Var(N)=\underbrace{\E[\Lambda]}_{\text{ruido de Poisson}}+\underbrace{\Var(\Lambda)}_{r\beta^2}=r\beta(1+\beta)>\E[N].
\]
La varianza supera a la media en exactamente la varianza de la tasa aleatoria; el \'indice de dispersi\'on es $\Var(N)/\E[N]=1+\beta$. Si la tasa fuera constante ($\Var\Lambda=0$) se recupera la Poisson, y $\text{BN}(r,\beta)\to\text{Poisson}(\lambda)$ cuando $\beta\to0$ y $r\to\infty$ con $r\beta=\lambda$.

\subsubsection*{Papel en la pr\'actica}
\begin{itemize}[leftmargin=*]
\item \textbf{Heterogeneidad no observada.} En una cartera de Auto u Hogar los asegurados no tienen todos la misma tasa de siniestros; $\Lambda$ es la tasa del asegurado, que no se conoce. Observamos recuentos mezclados, que son sobredispersos aunque cada asegurado sea Poisson.
\item \textbf{Correlaci\'on entre siniestros del mismo asegurado.} Un asegurado con muchos siniestros \emph{hoy} tiene una $\Lambda$ alta, y por tanto probablemente m\'as ma\~nana. Es la base matem\'atica de los sistemas bonus--malus y de la credibilidad.
\item \textbf{Factor com\'un de a\~no (sistem\'atico).} Un a\~no de tormentas eleva la tasa de toda la cartera; observado como un solo recuento, tambi\'en luce sobredisperso, y eso es riesgo no diversificable.
\item \textbf{Diagn\'ostico de segmentaci\'on.} La sobredispersi\'on residual, una vez incluidas las variables tarifarias, indica heterogeneidad que no se est\'a explicando (oportunidad de mejorar la tarificaci\'on).
\end{itemize}

\subsubsection*{Efecto en el capital}
Como $\Var(S)$ contiene $\Var(N)\E[X]^2$, ignorar la sobredispersi\'on subestima el riesgo. En el Caso~2 de Hogar, con $\Var(N)=600$ en vez de $150$ (Poisson), $\Var(S)$ pasa de $4{,}5\cdot10^8$ a $9{,}0\cdot10^8$ y el VaR Normal 99,5\,\% de $204.642$ a $227.275$ \euro\ (el modelo Poisson quedar\'ia un $10\%$ corto).

\paragraph{Conclusi\'on.} La BN es la distribuci\'on de un Poisson cuya tasa es aleatoria (Gamma): recoge la heterogeneidad y los factores comunes que producen una varianza mayor que la media, y esa varianza extra se paga en capital.

\bigskip\hrule\bigskip

\subsection{Pregunta 3 -- Ajuste MLE con Pareto de $\alpha=1{,}8$}

\noindent\textbf{Enunciado.} Si el ajuste MLE de severidad de una cartera arroja una Pareto con $\alpha=1{,}8$, ?`qu\'e implicaciones tiene para el c\'alculo del capital de solvencia y para la aproximaci\'on Normal Power?

\subsubsection*{Qu\'e momentos existen}
Para una Pareto (Lomax) de forma $\alpha$, $\E[X^k]<\infty$ si y solo si $k<\alpha$. Con $\alpha=1{,}8$:
\begin{itemize}[leftmargin=*,itemsep=1pt]
\item $\E[X]=\theta/(\alpha-1)=1{,}25\,\theta$ \textbf{existe} (se puede calcular la prima pura);
\item $\Var(X)=\infty$ y $\mu_3(X)=\infty$ (\textbf{no existen}).
\end{itemize}

\subsubsection*{Consecuencias}
\begin{enumerate}[leftmargin=*]
\item \textbf{La aproximaci\'on Normal no est\'a definida} (necesita $\Var(S)$) y \textbf{la Normal Power tampoco} (necesita $\mu_3$, que exige $\alpha>3$). No hay ``correcci\'on de asimetr\'ia'' que valga: ni siquiera tiene sentido el par\'ametro $\gamma_S$. Hay que usar Panjer, FFT o simulaci\'on, o bien teor\'ia de valores extremos.
\item \textbf{La varianza muestral no converge.} Al simular Pareto de $\alpha=1{,}8$ (media 1.000), la varianza muestral crece con el tama\~no de la muestra: $3{,}7\cdot10^6$ con $n=100$, $3{,}2\cdot10^6$ con 1.000, $9{,}1\cdot10^6$ con 10.000 y $7{,}2\cdot10^7$ con 100.000. Cualquier ``desviaci\'on t\'ipica'' estimada con datos es inestable y depende de la mayor observaci\'on.
\item \textbf{El capital lo fija un solo siniestro.} Para colas subexponenciales $P(S>x)\approx\E[N]\,P(X>x)$ (principio del ``gran salto \'unico''), de donde
\[
VaR_p(S)\approx F_X^{-1}\!\Big(1-\frac{1-p}{\E[N]}\Big)+\E[S]\qquad\text{(el resto de la cartera aporta su media).}
\]
Con $\E[N]=150$, $\E[X]=1.000$ (as\'i $\theta=800$) y $p=99{,}5\%$: $VaR\approx800\big[(0{,}005/150)^{-1/1{,}8}-1\big]+150.000=394.886$ \euro, unas $2{,}6$ veces la media. Esta f\'ormula reproduce el VaR exacto con un error del $0{,}2\%$ para $\alpha=1{,}5$ ($632.245$ frente a $633.750$), del $4{,}5\%$ para $\alpha=2$ y empeora cuando la cola es m\'as ligera ($16\%$ para $\alpha=3$).
\item \textbf{El capital es hipersensible a $\alpha$.} Con $\E[X]=1.000$ y frecuencia BN(50; 3), el VaR exacto 99,5\,\% sube de $249.750$ ($\alpha=3$) a $337.500$ ($\alpha=2$) y a $633.750$ ($\alpha=1{,}5$). Y $\alpha$ se estima con incertidumbre (error t\'ipico del MLE $\approx\hat\alpha/\sqrt n$): con 500 siniestros, $1{,}8\pm0{,}08$ ya cambia el capital de forma apreciable.
\item \textbf{El VaR deja de ser buena medida.} El TVaR existe ($\alpha>1$) pero es enorme respecto al VaR (con $\alpha=1{,}5$, $1{,}42$ M\euro\ frente a $0{,}63$ M\euro); el SST, basado en Expected Shortfall, exigir\'ia m\'as capital que Solvencia~II. Con $\alpha\le1$ ni la media existir\'ia.
\end{enumerate}

\subsubsection*{Qu\'e hacer}
\begin{itemize}[leftmargin=*,itemsep=1pt]
\item Contrastar el ajuste (gr\'afico $\log$--$\log$ de la cola, funci\'on de exceso medio lineal, estimador de Hill) y sustituir la Pareto global por una \emph{Pareto generalizada por encima de un umbral} (Tema~5).
\item \textbf{Acotar la cola}: l\'imites de p\'oliza y reaseguro de exceso de p\'erdida, de modo que la severidad retenida tenga todos los momentos.
\item Calcular el capital con Panjer/FFT o simulaci\'on sobre la severidad truncada, y a\~nadir un margen por incertidumbre de $\alpha$.
\item Documentar la elecci\'on y el an\'alisis de sensibilidad en la validaci\'on del modelo interno.
\end{itemize}

\paragraph{Conclusi\'on.} Con $\alpha=1{,}8$ el modelo ``media--varianza--asimetr\'ia'' deja de existir: la Normal y la Normal Power no son aplicables, el capital depende de un \'unico siniestro extremo y es muy sensible al par\'ametro. Se necesita c\'alculo num\'erico exacto o EVT, y reaseguro de cola.

\subsubsection*{Referencias}
\begin{itemize}[leftmargin=*,itemsep=0pt]
\item Klugman, S. A., Panjer, H. H. y Willmot, G. E. \emph{Loss Models: From Data to Decisions}. Wiley.
\item Kaas, R., Goovaerts, M., Dhaene, J. y Denuit, M. (2008). \emph{Modern Actuarial Risk Theory} (2.\textordfeminine\ ed.). Springer.
\item Panjer, H. H. (1981). Recursive evaluation of a family of compound distributions. \emph{ASTIN Bulletin}, 12.
\item Embrechts, P., Kl\"uppelberg, C. y Mikosch, T. (1997). \emph{Modelling Extremal Events}. Springer.
\end{itemize}

\end{document}
